% 自動生成（tools/mkchecklist.py p5）。各章の「ふりかえり」の項目を集めたもの（第13回）
\begingroup\footnotesize
\def\bkcb{\tikz[baseline=-0.5ex]\draw[\bkTF, line width=0.7pt, rounded corners=0.5pt] (0,-0.3em) rectangle (0.8em,0.5em);}
\begin{longtblr}{colspec={Q[c,m,wd=11mm] X[l,m] Q[c,m,wd=10mm] Q[c,m,wd=9mm] Q[c,m,wd=9mm]}, width=\linewidth, rowhead=1, colsep=3pt,
  row{1}={bg=bkVol!14, font=\sffamily\bfseries\scriptsize}, hlines={0.4pt,bkLine}, rowsep=0.7pt}
章 & できるようになったこと & ページ & 1回め & 2回め \\
{\sffamily\bfseries\color{\bkTT}原-0} & 原子1個の大きさと、1円玉の中の原子の数を、けたで言える。スマホのチップと原子炉が、原子のどの部分の物理で動くかを説明できる。 & \pageref{ph:原-0:sec3} & \bkcb & \bkcb \\
 & 何を無視したかを書いてから、原子の数・1個の大きさ・1個あたりのエネルギーを、7つの段で求められる。 & \pageref{ph:原-0:ex1} & \bkcb & \bkcb \\
 & 身近な「なぜ？」の答えが、この巻のどの章にあるかを探せる。 & \pageref{ph:原-0:sec4} & \bkcb & \bkcb \\
{\sffamily\bfseries\color{\bkTT}原-1} & 陰極線が電子の流れであることを、実験の結果から説明できる。 & \pageref{ph:原-1:sec1} & \bkcb & \bkcb \\
 & トムソンの方法で、電子の比電荷を求められる。 & \pageref{ph:原-1:sec2} & \bkcb & \bkcb \\
 & ミリカンの方法で、油滴の電荷と電気素量を求められる。 & \pageref{ph:原-1:sec3} & \bkcb & \bkcb \\
 & 原子核が小さく重いことを、いちばん近づく距離で説明できる。 & \pageref{ph:原-1:sec5} & \bkcb & \bkcb \\
 & 散乱角の式で、大きく曲がる α線の割合を計算できる。 & \pageref{ph:原-1:sec8} & \bkcb & \bkcb \\
 & 質量分析で、同位体のイオンが届く場所の違いを求められる。 & \pageref{ph:原-1:sec11} & \bkcb & \bkcb \\
{\sffamily\bfseries\color{\bkTT}原-2} & 熱した物の色が温度で変わるわけを、光の分布の山の動きで説明できる。 & \pageref{ph:原-2:sec1} & \bkcb & \bkcb \\
 & 光子のエネルギー $E=hf=hc/\lambda$ を、eV と J で求められる。 & \pageref{ph:原-2:sec3} & \bkcb & \bkcb \\
 & 光電効果のきまりを光子で説明し、$K_0=hf-W$ と阻止電圧を使える。 & \pageref{ph:原-2:sec5} & \bkcb & \bkcb \\
 & 光子の運動量 $p=h/\lambda$ から、放射圧を求められる。 & \pageref{ph:原-2:sec6} & \bkcb & \bkcb \\
 & コンプトン効果の式で、散乱した X線・γ線の波長とエネルギーを求められる。 & \pageref{ph:原-2:sec7} & \bkcb & \bkcb \\
 & モードの数と平均のエネルギーから、プランクの式・ウィーンの変位則・$\sigma$ を導ける。 & \pageref{ph:原-2:sec11} & \bkcb & \bkcb \\
{\sffamily\bfseries\color{\bkTT}原-3} & 加速した電圧から、電子の波長（ド・ブロイ波長）を求められる。 & \pageref{ph:原-3:sec2} & \bkcb & \bkcb \\
 & 電子の回折と、電子1個ずつの二重スリットのしま、電子顕微鏡が細かい物を見分けられるわけを、波の式で説明できる。 & \pageref{ph:原-3:sec5} & \bkcb & \bkcb \\
 & 波の集まり（波束）のはばから不確定性関係を導き、原子の中の電子のエネルギーの大きさを見積もれる。 & \pageref{ph:原-3:sec7} & \bkcb & \bkcb \\
 & 電子の波の分散関係から、群速度が粒の速さになることを示せる。 & \pageref{ph:原-3:sec9} & \bkcb & \bkcb \\
{\sffamily\bfseries\color{\bkTT}原-4} & 炎色反応と線スペクトルを、原子のエネルギーの段（準位）の差から説明できる。 & \pageref{ph:原-4:sec3} & \bkcb & \bkcb \\
 & ボーアの模型から水素原子のエネルギー準位を導き、スペクトルの波長を求められる。 & \pageref{ph:原-4:sec5} & \bkcb & \bkcb \\
 & X線管の電圧から最短波長を、モーズリーの法則から特性X線のエネルギーを求められる。 & \pageref{ph:原-4:sec9} & \bkcb & \bkcb \\
 & 誘導放出と反転分布、アインシュタインの係数から、レーザーのしくみを説明できる。 & \pageref{ph:原-4:sec11} & \bkcb & \bkcb \\
{\sffamily\bfseries\color{\bkTT}原-5} & 箱に閉じこめた電子のエネルギー $E_n=n^2h^2/(8mL^2)$ を、定在波と $\lambda=h/p$ から求められる。 & \pageref{ph:原-5:sec3} & \bkcb & \bkcb \\
 & $|\psi|^2$ を「見つかる確率」として読み、ある範囲に電子が見つかる確率を求められる。 & \pageref{ph:原-5:sec4} & \bkcb & \bkcb \\
 & シュレディンガー方程式を立て、無限に深い井戸を境界の条件で解ける。 & \pageref{ph:原-5:sec8} & \bkcb & \bkcb \\
 & 重ね合わせと期待値を計算し、差分法で固有値を数値で求められる。 & \pageref{ph:原-5:sec12} & \bkcb & \bkcb \\
{\sffamily\bfseries\color{\bkTT}原-6} & 壁の中で波動関数が $e^{-\kappa x}$ で弱まることを使い、すりぬける割合 $e^{-2\kappa a}$ を見積もれる。 & \pageref{ph:原-6:sec3} & \bkcb & \bkcb \\
 & フラッシュメモリや走査トンネル顕微鏡のしくみを、トンネル効果で説明できる。 & \pageref{ph:原-6:sec4} & \bkcb & \bkcb \\
 & 境目のつながりから、階段の反射率・有限の井戸の準位・正確な透過率を求められる。 & \pageref{ph:原-6:sec6} & \bkcb & \bkcb \\
 & 調和振動子の準位 $(n+\tfrac12)\hbar\omega$ を使い、分子の振動の光と比熱を説明できる。 & \pageref{ph:原-6:sec8} & \bkcb & \bkcb \\
{\sffamily\bfseries\color{\bkTT}原-7} & 周期表の横の並びの数（2・8・8・18…）が、原子の中の電子の「席」の数で決まることを説明できる。 & \pageref{ph:原-7:sec1} & \bkcb & \bkcb \\
 & 量子数で席を数えて $2n^2$ を導き、ナトリウム・鉄・銅などの電子配置を書ける。 & \pageref{ph:原-7:sec4} & \bkcb & \bkcb \\
 & 水素原子の電子の雲から、電子が見つかりやすい距離や見つかる確率を求め、電離エネルギーから有効核電荷を見積もれる。 & \pageref{ph:原-7:sec3} & \bkcb & \bkcb \\
 & 磁場の中の電子や陽子の向きが2つに分かれることから、ESR や MRI の周波数を求められる。 & \pageref{ph:原-7:sec6} & \bkcb & \bkcb \\
 & 3次元のシュレディンガー方程式に水素原子の $1s$ の波動関数を代入して $-13.6$ eV を確かめ、化学結合や分子の回転のしくみを読める。 & \pageref{ph:原-7:sec10} & \bkcb & \bkcb \\
{\sffamily\bfseries\color{\bkTT}原-8} & 金属・半導体・絶縁体の違いを、電子のエネルギーの帯（バンド）と、帯のすき間で説明できる。 & \pageref{ph:原-8:sec8} & \bkcb & \bkcb \\
 & ブラッグの条件で X線の回折の角を求め、格子定数から結晶の密度や原子の数を計算できる。 & \pageref{ph:原-8:sec2} & \bkcb & \bkcb \\
 & 3次元の箱で電子の席を数えて、金属のフェルミエネルギーとフェルミ速度を求められる。 & \pageref{ph:原-8:sec4} & \bkcb & \bkcb \\
 & フェルミ分布で、電子が席をうめる割合を求め、金属の電子の比熱が小さいわけを説明できる。 & \pageref{ph:原-8:sec5} & \bkcb & \bkcb \\
 & くし形の結晶の帯をプログラムで求め、平均自由行程・ウィーデマン・フランツの法則・デバイの $T^3$ 則を使える。 & \pageref{ph:原-8:sec10} & \bkcb & \bkcb \\
{\sffamily\bfseries\color{\bkTT}原-9} & 光電池に光を当てると電気ができ、LED は決まった向きにつないだときだけ光ることを、実験で確かめられる。 & \pageref{ph:原-9:sec1} & \bkcb & \bkcb \\
 & n形・p形の半導体の電子と正孔の数と抵抗率を求め、pn 接合が電流を一方の向きにだけ通すしくみを、図で説明できる。 & \pageref{ph:原-9:sec4} & \bkcb & \bkcb \\
 & ダイオードの電流と電圧の関係（約 60 mV で電流10倍）、LED の色と禁制帯の幅の関係を使える。 & \pageref{ph:原-9:sec7} & \bkcb & \bkcb \\
 & 太陽電池の電流と電圧のグラフから、開放電圧・いちばん大きな電力・曲線因子・効率を求められる。 & \pageref{ph:原-9:sec10} & \bkcb & \bkcb \\
 & フェルミ分布の裾からキャリアの数の式を、ドリフトと拡散のつり合いから内蔵電位とダイオードの式を導ける。 & \pageref{ph:原-9:sec11} & \bkcb & \bkcb \\
{\sffamily\bfseries\color{\bkTT}原-10} & スイッチ2つで AND と OR の回路を作り、オンとオフで2進数のたし算を表せる。 & \pageref{ph:原-10:sec1} & \bkcb & \bkcb \\
 & バイポーラトランジスタと MOSFET がスイッチ・増幅器としてはたらくしくみを pn 接合と電場で説明し、相互コンダクタンス・増幅度・オン抵抗を求められる。 & \pageref{ph:原-10:sec3} & \bkcb & \bkcb \\
 & CMOS の論理ゲートのしくみを説明し、チップの消費電力と信号のおくれを計算できる。 & \pageref{ph:原-10:sec7} & \bkcb & \bkcb \\
 & 理想オペアンプの反転・非反転増幅の倍率を求め、センサーの小さな電圧を A/D 変換する流れを説明できる。 & \pageref{ph:原-10:sec8} & \bkcb & \bkcb \\
 & MOSFET の電流の式を反転層の電荷とドリフトから導き、パワー半導体の損失と、ニュートン法で回路を解く方法を使える。 & \pageref{ph:原-10:sec10} & \bkcb & \bkcb \\
{\sffamily\bfseries\color{\bkTT}原-11} & 原子番号・質量数・同位体ということばで、原子核のなかみを説明できる。 & \pageref{ph:原-11:sec3} & \bkcb & \bkcb \\
 & 原子の質量の表から、結合エネルギーと核反応の Q 値を計算できる。 & \pageref{ph:原-11:sec4} & \bkcb & \bkcb \\
 & 光の時計から時間のおくれを導き、GPS の時計のずれを計算できる。 & \pageref{ph:原-11:sec7} & \bkcb & \bkcb \\
 & $E^2=(pc)^2+(mc^2)^2$ から、$E=mc^2$ とコンプトンの式を導ける。 & \pageref{ph:原-11:sec8} & \bkcb & \bkcb \\
 & 半経験的質量公式で、結合エネルギーの曲線の形を説明できる。 & \pageref{ph:原-11:sec10} & \bkcb & \bkcb \\
{\sffamily\bfseries\color{\bkTT}原-12} & α崩壊・β崩壊・γ線で、原子核がどう変わるかを書ける。 & \pageref{ph:原-12:sec2} & \bkcb & \bkcb \\
 & 半減期を使って、残る量や、たった時間を求められる。 & \pageref{ph:原-12:sec3} & \bkcb & \bkcb \\
 & 放射能（Bq）と原子の数の関係 $A=\lambda N$ を使える。 & \pageref{ph:原-12:sec4} & \bkcb & \bkcb \\
 & 崩壊の法則を微分方程式で書き、数えるばらつきを見積もれる。 & \pageref{ph:原-12:sec5} & \bkcb & \bkcb \\
 & 親と子の崩壊（放射平衡）と、α崩壊のトンネルを式で読める。 & \pageref{ph:原-12:sec8} & \bkcb & \bkcb \\
{\sffamily\bfseries\color{\bkTT}原-13} & α線・β線・γ線（X線）・中性子線が、どんな物で止まるか（透過力）と、物を電離するしくみを説明できる。 & \pageref{ph:原-13:sec1} & \bkcb & \bkcb \\
 & 光子の減弱の式 $I=I_0e^{-\mu x}$ と半価層を使って、X線写真や遮へいの厚さを計算できる。 & \pageref{ph:原-13:sec3} & \bkcb & \bkcb \\
 & 光電効果・コンプトン散乱・電子対生成の違いと、荷電粒子・中性子の止まり方の違いを説明できる。 & \pageref{ph:原-13:sec4} & \bkcb & \bkcb \\
 & 阻止能とブラッグピーク、断面積と平均自由行程、中性子の減速を、式で読める。 & \pageref{ph:原-13:sec7} & \bkcb & \bkcb \\
 & 放射化の式 $A=N\sigma\phi(1-e^{-\lambda t})$ で、飽和放射能に近づくようすを計算できる。 & \pageref{ph:原-13:sec9} & \bkcb & \bkcb \\
{\sffamily\bfseries\color{\bkTT}原-14} & 放射線は目に見えないが、電離を使って数えられることと、Bq・Gy・Sv の意味の違いを説明できる。 & \pageref{ph:原-14:sec2} & \bkcb & \bkcb \\
 & 吸収線量から等価線量・実効線量を計算し、自然放射線や医療の線量と比べられる。 & \pageref{ph:原-14:sec5} & \bkcb & \bkcb \\
 & 点線源の逆2乗の法則と遮へいを組み合わせ、時間・距離・遮へいの3原則で線量を見積もれる。 & \pageref{ph:原-14:sec7} & \bkcb & \bkcb \\
 & 計数のばらつき（$\sqrt N$）・数え落とし・検出限界・立体角を使って、測定の結果を正しく読める。 & \pageref{ph:原-14:sec3} & \bkcb & \bkcb \\
 & 確定的影響と確率的影響の違い、線量限度の意味、低い線量のリスクについての2つの見方を説明できる。 & \pageref{ph:原-14:sec11} & \bkcb & \bkcb \\
{\sffamily\bfseries\color{\bkTT}原-15} & ウラン235 の核分裂で、1回に約 200 MeV のエネルギーと2〜3個の中性子が出ることと、連鎖反応・原子炉のしくみを、図で説明できる。 & \pageref{ph:原-15:sec1} & \bkcb & \bkcb \\
 & 増倍率 $k$・反応度・遅発中性子を使って、原子炉の出力の増え方（炉の周期）を見積もれる。 & \pageref{ph:原-15:sec6} & \bkcb & \bkcb \\
 & 4因子公式と1群の拡散方程式で、増倍率と、原子炉が臨界になる大きさを求められる。 & \pageref{ph:原-15:sec7} & \bkcb & \bkcb \\
 & 1点炉の動特性の式・反応度係数・キセノンの式を、微分方程式として読める。 & \pageref{ph:原-15:sec8} & \bkcb & \bkcb \\
 & 崩壊熱と核融合のエネルギーを数で見積もり、原子炉を止めたあとも冷やし続けなければならないわけを説明できる。 & \pageref{ph:原-15:sec11} & \bkcb & \bkcb \\
{\sffamily\bfseries\color{\bkTT}原-16} & 原子の力の使いみち（発電・医療・工業・農業・研究）と、核分裂の発見から原子炉・原子爆弾・原子力発電までの歴史を説明できる。 & \pageref{ph:原-16:sec1} & \bkcb & \bkcb \\
 & 使い終わった燃料が熱と放射線を出し続けるわけを、崩壊熱と半減期で説明し、冷やす期間と、人から離しておく期間の桁を見積もれる。 & \pageref{ph:原-16:ex2} & \bkcb & \bkcb \\
 & 福島第一原子力発電所の事故の経過と、廃炉・処理水のことを、物理の量（崩壊熱・Bq・Sv）と公的な資料の数で説明できる。 & \pageref{ph:原-16:sec5} & \bkcb & \bkcb \\
 & 放射線や原子力の数を、何を・どの基準で数えた値かに気をつけて読み、物理でわかることと、社会が選ぶことを分けて考えられる。 & \pageref{ph:原-16:sec9} & \bkcb & \bkcb \\
\end{longtblr}
\endgroup
